\section{Introduction}
\label{sec:intro}



Humans naturally integrate information from multiple modalities (e.g., vision, language, and audio) to make intelligent decisions~\cite{baltruvsaitis2018multimodal}. Consequently, multimodal learning has emerged as a core paradigm in the research community, aiming to leverage multimodal data to enhance decision-making across diverse application domains, including healthcare assistance~\cite{ghosh2024clipsyntel} and malicious content detection~\cite{kiela2020hateful}. However, in open-world environments, multimodal models often struggle to maintain performance under distribution shifts~\cite{liang2025comprehensive}, which frequently arise due to unexpected environmental changes or noise, such as weather variations and data corruption. To address this challenge, Test-Time Adaptation (TTA) has gained increasing attention in multimodal settings. TTA enables pre-trained models to adapt to unseen target domains at inference time without access to source data, thereby reducing computational overhead by eliminating the need for retraining and preserving source data privacy. Recently, a number of promising TTA methods have shown great performance in open-world scenarios.

Real-world scenarios, however, present an even more challenging problem: continuously changing target domains \cite{wang2022continual}. Specifically, in the senarios of autonomous driving, its system must rapidly adapt to vary weather conditions, lighting changes, and road conditions.